\chapter{Manual de instalación}

En este capítulo se detalla el procedimiento para instalar y ejecutar la aplicación. Gracias al uso
de contenedores, el despliegue es reproducible y no requiere instalar manualmente las dependencias
de Python ni de Node.js: basta con disponer de Docker.

El proyecto contempla dos formas de ejecución, pensadas para propósitos distintos. La
\textbf{ejecución en local} (secciones~\ref{sec:requisitos} y~\ref{sec:lanzamiento}) levanta los dos
servicios por separado con Docker Compose y recarga en caliente, que es lo adecuado durante el
desarrollo. El \textbf{despliegue en un servidor público} (sección~\ref{sec:despliegue}) emplea, en
cambio, una única imagen autocontenida pensada para producción.

\section{Requisitos previos}
\label{sec:requisitos}

Para ejecutar la aplicación es necesario contar con:

\begin{itemize}
	\item \textbf{Docker}~\cite{docker} y \textbf{Docker Compose} instalados y en ejecución. En Windows y macOS
		se recomienda instalar \textit{Docker Desktop}; en Linux, el motor de Docker junto con el
		\textit{plugin} de Compose.
	\item Los \textbf{ficheros de datos} del atlas de ejemplo (\texttt{AAL90.node} y
		\texttt{AAL90.edge}) ubicados en el directorio \texttt{backend/data/} del proyecto.
	\item Un \textbf{navegador moderno} con soporte de WebGL 2.0 (Chrome, Firefox o Edge).
\end{itemize}

Puede comprobarse que Docker está correctamente instalado con:

\begin{verbatim}
docker --version
docker compose version
\end{verbatim}

\section{Instalación y lanzamiento en local}
\label{sec:lanzamiento}

Una vez clonado o descargado el repositorio, se accede a su directorio raíz y se construyen y
levantan los contenedores.

\subsection{Alternativa 1: recomendada}

Se emplea el \textit{script} auxiliar incluido en el proyecto, que levanta los servicios y muestra
por consola los puertos asignados automáticamente:

\begin{verbatim}
./up-brain-viz.sh
\end{verbatim}

\subsection{Alternativa 2}

De forma equivalente, se pueden construir y levantar los contenedores directamente con Docker
Compose:

\begin{verbatim}
docker compose up --build
\end{verbatim}

\subsection{Post-ejecución y resumen}

Los servicios de \textit{backend} y \textit{frontend} exponen sus puertos internos (5000 y 5173,
respectivamente) en puertos del anfitrión asignados automáticamente, evitando conflictos con otros
procesos. Para conocer los puertos publicados se ejecuta:

\begin{verbatim}
docker compose ps
\end{verbatim}

A partir de la columna de puertos publicados se obtiene la dirección de acceso; por ejemplo, el
frontend estará disponible en \texttt{http://localhost:<puerto>}. Para detener la aplicación:

\begin{verbatim}
docker compose down
\end{verbatim}

\subsection{Consideraciones}

\begin{itemize}
	\item El navegador se comunica únicamente con el \textit{frontend}; las peticiones a
		\texttt{/api} y \texttt{/health} se redirigen internamente al \textit{backend} dentro de la
		red de contenedores.
	\item El entorno está configurado para desarrollo con recarga en caliente
		(\textit{hot-reload}): los cambios en el código fuente se reflejan sin reconstruir la imagen.
	\item El \textit{script} \texttt{get-docker.sh} incluido en el repositorio es un instalador de
		Docker para Linux; en macOS y Windows debe usarse Docker Desktop.
\end{itemize}

\section{Despliegue en un servidor público}
\label{sec:despliegue}

Además de la ejecución en local, la aplicación se despliega en un servidor accesible desde Internet
para que el director del proyecto pueda validar cada \textit{sprint} y para que el tribunal pueda
probarla sin necesidad de instalar nada.

\subsection{Imagen de producción}

El modo de desarrollo descrito en la sección anterior no es adecuado para un servidor público por dos
motivos: emplea servidores de desarrollo ---tanto el de Vite como el integrado de Flask, que no están
pensados para producción y no deben ejecutarse con el modo de depuración activo en una dirección
pública--- y reparte la aplicación en dos servicios que se comunican mediante un nombre interno de la
red de Docker Compose, inexistente fuera de ella.

Por ello se ha añadido un \texttt{Dockerfile} en la raíz del proyecto que genera una \textbf{imagen
única de producción} en dos etapas:

\begin{enumerate}
	\item En la primera etapa se compila la aplicación de React, obteniendo un conjunto de ficheros
		estáticos optimizados.
	\item En la segunda, esos ficheros se copian a la imagen de Python, donde el propio servidor sirve
		tanto la interfaz como la API.
\end{enumerate}

De este modo todo el sistema queda en un solo contenedor, desaparece la necesidad de un \textit{proxy}
entre servicios y la imagen puede desplegarse en cualquier plataforma que admita contenedores. La
aplicación se ejecuta con \textbf{gunicorn}, un servidor WSGI de producción, escuchando en todas las
interfaces y en el puerto que indique la variable de entorno \texttt{PORT}, tal y como esperan las
plataformas de despliegue habituales. La imagen puede construirse y probarse en local con:

\begin{verbatim}
docker build -t brain-viz:prod .
docker run -p 8080:8080 brain-viz:prod
\end{verbatim}

\subsection{Publicación}

El repositorio incluye un fichero \texttt{render.yaml} que describe el servicio, de forma que la
plataforma de despliegue detecta la configuración automáticamente a partir del repositorio: tipo de
servicio, ruta del \texttt{Dockerfile}, región y ruta de comprobación de estado (\texttt{/health}).
Cada nueva versión publicada en la rama principal del repositorio desencadena un nuevo despliegue.

% NOTA: insertar aquí la URL pública una vez completado el despliegue.

\subsection{Consideraciones del entorno gratuito}

El servicio se aloja en un plan gratuito, lo que conlleva una limitación que conviene conocer: la
instancia se suspende tras un periodo de inactividad, de modo que la primera petición tras ese
periodo puede tardar aproximadamente un minuto mientras el servicio se reactiva; las siguientes se
atienden con normalidad. Para reducir este efecto durante los periodos de evaluación se emplea un
servicio externo de comprobación periódica que consulta el \textit{endpoint} \texttt{/health} a
intervalos regulares, manteniendo la instancia activa.
